\chapter{Osservabilit\`a e monitoring}
\label{chap:observability}

\section{Strategia di osservabilit\`a}
\label{sec:observability_strategy}

La piattaforma espone tre canali di telemetria complementari, che coprono livelli
differenti di granularit\`a:

\begin{enumerate}[leftmargin=*]
    \item \textbf{Log strutturati} (livello infrastrutturale): emessi in stdout da ogni
          container, aggregabili da qualsiasi stack di log management.
    \item \textbf{Health endpoint} (livello servizio): probe attivo per Kubernetes.
    \item \textbf{Agent operations log} (livello applicativo/agentico): audit trail
          persistente delle decisioni dell'agente.
\end{enumerate}

\section{Log strutturati}
\label{sec:logging}

Tutti i microservizi emettono log in stdout con il formato:
\begin{verbatim}
%(asctime)s [%(levelname)s] %(name)s: %(message)s
\end{verbatim}

Questo formato \`e conforme alla convenzione \emph{12-factor app}: i log sono trattati
come stream di eventi e possono essere raccolti da qualsiasi sistema di aggregazione
(Loki, Elasticsearch, CloudWatch, ecc.) senza configurazione applicativa.

Gli eventi significativi loggati includono:

\begin{itemize}[leftmargin=*]
    \item \textbf{email\_poller}: inizio e fine di ogni ciclo di polling, numero di
          account elaborati, UID massimo ricevuto, eccezioni per account.
    \item \textbf{agent\_worker}: avvio del loop, job/file presi in carico, risultati
          di Tika, nome dell'indice OpenSearch creato, decisioni dell'agente,
          errori con stacktrace completo.
    \item \textbf{mcp\_server}: ogni tool call con \texttt{user\_id}, \texttt{parent\_id}
          e numero di risultati restituiti.
    \item \textbf{backend}: richieste HTTP (via Uvicorn access log), errori di
          validazione, eccezioni non gestite.
\end{itemize}

\section{Health check endpoint}
\label{sec:health}

L'endpoint \texttt{GET /health} esegue due verifiche attive:

\begin{enumerate}[leftmargin=*]
    \item \textbf{Database}: \texttt{SELECT 1} tramite la sessione SQLAlchemy.
          Fallimento $\rightarrow$ \texttt{status: "unhealthy"}, HTTP 503.
    \item \textbf{S3}: creazione del client boto3.
          Fallimento $\rightarrow$ \texttt{s3: "error"} ma \texttt{status: "healthy"},
          HTTP 200. Questa asimmetria riflette l'impatto reale: il read path dei
          metadati rimane funzionante anche senza S3.
\end{enumerate}

L'endpoint \`e progettato per essere utilizzato direttamente come:
\begin{itemize}[leftmargin=*]
    \item \textbf{Kubernetes liveness probe}: il container viene riavviato se restituisce
          503 continuativamente.
    \item \textbf{Kubernetes readiness probe}: il container viene rimosso dal load
          balancer durante la fase di startup o di recovery del database.
\end{itemize}

\section{Agent Operations Log (audit trail agentico)}
\label{sec:audit_log}

La tabella \texttt{agent\_operations} costituisce il registro immutabile di ogni azione
intrapresa dall'agente. \`E esposta all'endpoint \texttt{GET /agent/operations} e
visualizzata nella \emph{Agent History view} del frontend.

Ogni record contiene:
\begin{itemize}[leftmargin=*]
    \item \texttt{operation\_type}: tipo di evento in un insieme chiuso
          (\texttt{email\_received}, \texttt{attachment\_extracted}, \texttt{auto\_filed},
          \texttt{sent\_to\_inbox}, \texttt{duplicate\_skipped},
          \texttt{manual\_upload\_received}, \texttt{error}).
    \item \texttt{description}: testo human-readable dell'operazione.
    \item \texttt{details} (JSONB): metadati strutturati specifici per tipo di evento
          (es.\ \texttt{folder\_id}, \texttt{confidence}, \texttt{sender}).
    \item \texttt{job\_id} e \texttt{agent\_file\_id}: riferimenti che permettono
          la tracciabilit\`a completa dall'email originale alla decisione finale.
\end{itemize}

La presenza dell'\texttt{eml\_storage\_key} nel job referenziato permette il
\textbf{replay forensico}: a partire dal file \texttt{.eml} originale \`e sempre
possibile ricostruire il contesto in cui una decisione \`e stata presa, soddisfacendo
il requisito di auditabilit\`a delle linee guida.

\section{Frontend Transparency Layer}
\label{sec:transparency}

La \emph{Agent History view} sul frontend costituisce il \emph{transparency plane}: una
timeline continua che distingue visivamente le azioni umane (\texttt{history}) da quelle
agentiche (\texttt{agent\_operations}). I badge sulle cartelle modificate dall'agente
offrono una notifica leggera senza richiedere un canale push dedicato.

La \emph{Inbox view} espone le proposte in attesa con il reasoning testuale dell'agente
(\texttt{agent\_reasoning}), rendendo trasparente all'utente il processo decisionale
del sistema.

\section{Metriche e dashboard}
\label{sec:metrics}

\TODO{Questa sezione descriver\`a l'integrazione con Prometheus per la raccolta di
metriche operative (latenza per endpoint, throughput di job/min, tasso di auto-filing,
tasso di errori per componente) e con Grafana per la visualizzazione. Includer\`a
dashboard per: (1) salute dei microservizi, (2) pipeline agentica (job in pending/processing,
classificazioni per ora, distribuzione delle confidence), (3) performance OpenSearch
(latenza query, dimensione indici per tenant).}

\section{Distributed tracing}
\label{sec:tracing}

\TODO{Integrazione OpenTelemetry per il tracciamento delle richieste attraverso i
microservizi. Gli span tracceranno il percorso completo da email ricevuta a documento
archiviato, con attributi per ogni step della pipeline (Tika latency, LLM latency,
OpenSearch indexing latency). L'exporter sar\`a configurato verso Jaeger o Tempo.}
